\noindent
This manual documents, in a self-contained way, the complete mathematics of the
fitting engine for $^{57}$Fe Mössbauer spectra implemented in the application. It
covers (i)~the \emph{forward model} —line profiles, positions and intensities of the
six lines of the sextet with first-order and Kündig treatments of the
magnetic–quadrupole coupling, doublets and singlets—; (ii)~the \emph{discrete fit}
—data statistics in Gauss and Poisson modes, propagation of the calibration
uncertainty, cost functional with robust losses, trust-region optimization with
multistart and a global option, covariance, diagnostics and bootstrap—; and
(iii)~the \emph{distribution fit} $P(\BHF)$ or $P(\dEQ)$ —discretization, kernel,
Tikhonov and total-variation regularization, conditioning, effective degrees of
freedom, automatic selection of the regularization parameter, error bars and
parametric distributions. References to the file and function of the code are
included.



\bigskip

%======================================================================
\section{Forward model}
%======================================================================

\subsection{Transmission spectrum}\label{sec:fwd}

After folding the raw data around the fold point, the spectrum is defined over $N$
channels with velocities $\{v_i\}$ and normalized transmission $\{y_i\}$, with
$y_i\approx 1$ away from the absorption peaks. The forward model is a linear
superposition of absorptions on top of a baseline:
\begin{equation}
  \boxed{\;
  T(v;\bm\theta) \;=\; b + s\,v \;-\; \sum_{c=1}^{N_c} A_c(v;\bm\theta_c)
  \;}
  \label{eq:Tmodel}
\end{equation}
where $b$ is the baseline (multiplicative; the data are normalized by the factor
\code{norm\_factor}), $s$ the slope, and each $A_c$ is the absorption of a
component. Besides sextet, doublet and singlet, the components can be
\textbf{magnetic relaxation} models (phenomenological relaxation, Blume--Tjon,
Néel; see the relaxation document). Implemented in \file{core/physics.py},
function \code{total\_model}.

By default it is a \emph{thin-absorber} model (linear in depth). Optionally a
\textbf{thick-absorber} model with saturation is applied,
$T = b + s\,v - C\bigl(1-e^{-A_{\text{tot}}/C}\bigr)$ (parameter $C=$
\code{sat\_scale}); the derivation is in the thickness-correction document.

\subsubsection{Drive waveform and velocity axis}\label{sec:drive}
With a \textbf{constant-acceleration} drive (triangular wave) the velocity is
linear in the channel: after folding, $\{v_i\}$ is \textbf{equispaced} on
$[-v_{\max},v_{\max}]$ (\code{velocity\_axis} in \code{core/folding.py}). With a
\textbf{sinusoidal} drive the velocity is not linear; following the NORMOS
criterion (\code{TRIANG=.FALSE.}, \code{FOLD=.FALSE.}, \code{SIMULT=.TRUE.}) the
data are \emph{not folded} and each channel is assigned its actual velocity
\begin{equation}
  v_i \;=\; v_{\max}\,\sin\!\Big(\frac{2\pi\,(i-c_0)}{N}\Big),
\end{equation}
with $c_0$ the phase (channel of $v=0$), auto-detected by symmetry and refinable as
a fit parameter. The resulting axis is not monotonic, which the model tolerates
because it is evaluated point by point (\code{sine\_velocity\_axis} in \code{core/folding.py}).

\subsection{Line profiles}\label{sec:lineshape}

Each elementary line uses a profile normalized to unit height at its center. The
application offers two profiles, selectable globally.

\subsubsection{Lorentzian}
\begin{equation}
  L(v;v_0,\gamma) \;=\; \frac{(\gamma/2)^{2}}{(v-v_0)^{2}+(\gamma/2)^{2}},
  \label{eq:lorentz}
\end{equation}
where $\gamma$ is the \textbf{width parameter that is fitted}: the \emph{full width
at half maximum} (FWHM). The half-width (HWHM) is therefore $\gamma/2$, which is
what appears in the formula. It satisfies $L(v_0)=1$.

\begin{nota}[Width convention]
In the code and in the interface, the width parameter ($\gamma_1$, etc.) is the
\textbf{FWHM}. Internally the profile uses the half-width $\gamma/2$
(\code{lorentzian} in \code{core/physics.py} divides by two before evaluating).
Reading $\gamma$ as HWHM (as in old versions of this document) gave widths off by a
factor of 2.
\end{nota}

\subsubsection{Pseudo-Voigt (Faddeeva)}
The convolution of a Lorentzian (HWHM $\gamma/2$) with a Gaussian (width $\sigma$)
is the Voigt profile, evaluated through the Faddeeva function
$w(z)=e^{-z^2}\erfc(-iz)$:
\begin{equation}
  \widetilde V(v;v_0,\gamma,\sigma) \;=\;
   \frac{1}{\sigma\sqrt{2\pi}}\,\Re\!\left[w\!\left(\frac{(v-v_0)+i\gamma/2}{\sigma\sqrt2}\right)\right].
\end{equation}
So that the profile has height exactly $1$ at $v_0$ \emph{independently of the
sampling}, it is normalized by its analytical value at the peak:
\begin{equation}
  V(v;v_0,\gamma,\sigma) \;=\; \frac{\widetilde V(v;v_0,\gamma,\sigma)}{\widetilde V(v_0;v_0,\gamma,\sigma)},
  \qquad
  \widetilde V(v_0) = \frac{1}{\sigma\sqrt{2\pi}}\,\Re\!\left[w\!\left(\tfrac{i\gamma/2}{\sigma\sqrt2}\right)\right].
  \label{eq:voigt}
\end{equation}
Since $w(iy)=e^{y^2}\erfc(y)\in\mathbb R^{+}$ for $y\in\mathbb R$, the denominator
is real and positive. In the limit $\sigma\to0$, $V\to L$.

\begin{nota}[Normalization to the peak]
Normalizing by the \emph{discrete} maximum on the grid (as previous versions did)
underestimates the true peak when no $v_i$ falls at $v_0$, inflating the amplitude
and making the depth depend on the sampling. Form \eqref{eq:voigt} removes that
bias.
\end{nota}

\begin{figure}[htbp]\centering
  \includegraphics[width=0.95\textwidth]{apx_lineshapes}
  \caption{Line profiles normalized to unit height and to equal FWHM. The
  Lorentzian \eqref{eq:lorentz} has \emph{heavy tails}; the Gaussian decays much
  faster; the pseudo-Voigt \eqref{eq:voigt} interpolates between the two. The
  logarithmic scale (right) makes the weight of the tails visible, which is what
  distinguishes one profile from another far from the center.}
  \label{fig:apx_lineshapes}
\end{figure}

\subsection{Magnetic sextet}\label{sec:sextet}

A sextet arises from the nuclear Zeeman splitting combined, where applicable, with
the quadrupole interaction. The six lines correspond to the allowed transitions
($\Delta m=0,\pm1$) between the ground state $I_g=\tfrac12$ and the excited state
$I_e=\tfrac32$. The absorption is
\begin{equation}
  A_{\text{sext}}(v) \;=\; d\,\sum_{j=1}^{6} W_j\, P\!\left(v;v_{0,j},\gamma_j\right),
  \label{eq:Asext}
\end{equation}
with $d$ the depth, $P$ the profile (\S\ref{sec:lineshape}), $W_j$ the intensities
(\S\ref{sec:intensities}) and $\gamma_j$ the widths
\begin{equation}
  \bm\gamma = \gamma_1\,[\,1,\;\gamma_2,\;\gamma_3,\;\gamma_3,\;\gamma_2,\;1\,],
\end{equation}
where $\gamma_2,\gamma_3$ are \emph{relative} widths. The positions $v_{0,j}$ are
obtained by one of two treatments. Figure~\ref{fig:apx_zeeman} summarizes the
physical origin of the six lines.

\begin{figure}[htbp]\centering
  \includegraphics[width=0.82\textwidth]{apx_zeeman}
  \caption{Origin of the sextet. The hyperfine field Zeeman-splits the ground level
  ($I_g=\tfrac12$, two sublevels) and the excited level ($I_e=\tfrac32$, four
  sublevels). The six allowed transitions ($\Delta m=0,\pm1$) produce the six lines,
  with intensities $3:2:1:1:2:3$ in a powder and positions at $33$~T of
  $\pm0{,}839/\pm3{,}084/\pm5{,}329$~mm/s. The quadrupole (not shown) rigidly shifts
  this pattern according to \eqref{eq:sext_pos1}.}
  \label{fig:apx_zeeman}
\end{figure}

\subsubsection{First-order treatment}
The quadrupole is treated as a rigid perturbation on the Zeeman pattern:
\begin{equation}
  v_{0,j} \;=\; \frac{\BHF}{B_0}\,r_j^{(0)} \;+\; \delta \;+\; s_j\,\dEQ,
  \qquad j=1,\dots,6,
  \label{eq:sext_pos1}
\end{equation}
with $B_0=33\,$T the reference field, $\delta$ the isomer shift, $r^{(0)}$ the
positions calibrated at $B_0$ (Appendix~\ref{ap:calib}) and the quadrupole pattern
\begin{equation}
  \bm s = \left[\,+\tfrac12,\,-\tfrac12,\,-\tfrac12,\,-\tfrac12,\,-\tfrac12,\,+\tfrac12\,\right].
\end{equation}
Valid when $\dEQ \ll \mu_N g_e \BHF$ and the axis of the electric field gradient
(EFG) is parallel to $\vec B$.

\begin{figure}[htbp]\centering
  \includegraphics[width=0.9\textwidth]{apx_quadrupole}
  \caption{First-order effect of the quadrupole \eqref{eq:sext_pos1}. Compared with
  the pure sextet ($\Delta E_Q=0$, gray), a $\Delta E_Q\neq0$ shifts the lines
  according to the pattern $\bm s=[+\tfrac12,-\tfrac12,-\tfrac12,-\tfrac12,-\tfrac12,+\tfrac12]$:
  the two outer lines move one way and the four inner lines the other, breaking the
  symmetry of the sextet. The magnitude of the shift measures $\Delta E_Q$.}
  \label{fig:apx_quadrupole}
\end{figure}

\subsubsection{Kündig treatment (full axial Hamiltonian)}\label{sec:kundig}
For an axial EFG (asymmetry parameter $\eta=0$) and arbitrary angle $\beta$ between
$\vec B$ and the principal axis $V_{zz}$, the positions are obtained by
diagonalizing the excited-state Hamiltonian. In the basis
$\{\lvert\tfrac32\rangle, \lvert\tfrac12\rangle, \lvert{-}\tfrac12\rangle,
\lvert{-}\tfrac32\rangle\}$:
\begin{equation}
  \mathcal H_e \;=\; \omega_e\, I_z \;+\; \frac{\dEQ}{6}\bigl(3 I_{z'}^2 - I(I+1)\bigr),
  \qquad I_{z'} = I_z\cos\beta + I_x\sin\beta,
  \label{eq:Hkundig}
\end{equation}
with $I(I+1)=\tfrac{15}{4}$ and $\omega_e=(\BHF/B_0)\,\omega_e^{(0)}$ the Zeeman
splitting of the excited state. Numerical diagonalization of the $4\times4$ matrix
gives the eigenvalues $E^{e}_{m}$ and eigenvectors. The ground state
($I_g=\tfrac12$) is pure Zeeman:
\begin{equation}
  E^{g}_{\pm 1/2} = \pm\tfrac12\,\omega_g, \qquad \omega_g=(\BHF/B_0)\,\omega_g^{(0)}.
\end{equation}
The six positions are the energy differences of the allowed transitions, assigning
to each excited-state eigenvalue its dominant $m_e$ (by eigenvector overlap):
\begin{equation}
\begin{aligned}
  v_{0,1}&=E^{e}_{-3/2}-E^{g}_{-1/2}, & v_{0,2}&=E^{e}_{-1/2}-E^{g}_{-1/2}, & v_{0,3}&=E^{e}_{+1/2}-E^{g}_{-1/2},\\
  v_{0,4}&=E^{e}_{-1/2}-E^{g}_{+1/2}, & v_{0,5}&=E^{e}_{+1/2}-E^{g}_{+1/2}, & v_{0,6}&=E^{e}_{+3/2}-E^{g}_{+1/2},
\end{aligned}
\end{equation}
plus $\delta$. In the limit $\beta=0$ it reproduces \eqref{eq:sext_pos1} to machine
precision. Implemented in \file{core/hamiltonian.py}.

\subsubsection{Polycrystalline average}\label{sec:powder}
In a powder sample all orientations $\beta$ are equiprobable. The absorption is
averaged over the sphere:
\begin{equation}
  A_{\text{sext}}^{\text{pol}}(v) \;=\; \frac12\int_0^{\pi}\!\sin\beta\;A_{\text{sext}}(v;\beta)\,d\beta
  \;\approx\; \sum_{k=1}^{N_q} w_k\, A_{\text{sext}}(v;\beta_k),
  \label{eq:powder}
\end{equation}
with Gauss–Legendre quadrature nodes and weights over $x=\cos\beta\in[-1,1]$,
$\beta_k=\arccos x_k$, $w_k=\tfrac12 w_k^{\mathrm{GL}}$ (so that
$\sum_k w_k=1$). The computation is vectorized: for $N_q$ orientations a tensor
$(N_q,4,4)$ is built and diagonalized at once. Default value $N_q=20$ (error
$<0{,}01\%$).

\begin{figure}[htbp]\centering
  \includegraphics[width=0.9\textwidth]{apx_powder}
  \caption{Polycrystalline average when $\Delta E_Q\neq0$. In a single crystal the
  positions of the six lines depend on the angle $\beta$ between $\vec B$ and the
  EFG axis (curves $\beta=0^\circ$ and $\beta=90^\circ$, computed with the Kündig
  Hamiltonian \eqref{eq:Hkundig}). In a powder, all orientations coexist: the
  spectrum is the weighted sum \eqref{eq:powder} over $\beta$, which
  \emph{broadens} and deforms the lines (upper curve). This is why the correct
  treatment of a quadrupolar sextet in a powder requires averaging, not a single
  orientation.}
  \label{fig:apx_powder}
\end{figure}

\subsection{Intensities and texture}\label{sec:intensities}

The intensities of the six lines respect the symmetry $1\!\leftrightarrow\!6$,
$2\!\leftrightarrow\!5$, $3\!\leftrightarrow\!4$:
\begin{equation}
  \bm W = [\,I_1,\,I_2,\,I_3,\,I_3,\,I_2,\,I_1\,],\qquad
  I_1=i_3 i_1,\;\; I_2=i_3 i_2,\;\; I_3=i_3.
\end{equation}
The application offers two parametrizations per sextet.

\subsubsection{Free intensities}
$i_1,i_2,i_3$ are independent fit parameters. It does not impose the $3{:}2{:}1$
ratio; useful for numerical flexibility and empirical interpretation.

\subsubsection{Texture parameter}\label{sec:texture}
For a sample with angle $\theta$ between the magnetic moment and the direction of
the $\gamma$ photon, the magnetic-dipole amplitudes are
\begin{equation}
  W_{1,6}\propto 3(1+\cos^2\theta),\quad W_{2,5}\propto 4\sin^2\theta,\quad W_{3,4}\propto (1+\cos^2\theta).
\end{equation}
Defining the texture parameter $t=\sin^2\theta\in[0,1]$ and normalizing by
$(1+\cos^2\theta)=2-t$:
\begin{equation}
  \boxed{\;i_1 = 3,\qquad i_2 = \frac{4t}{2-t},\qquad i_3 = 1.\;}
  \label{eq:texture}
\end{equation}
Cases: $t=\tfrac23\Rightarrow 3{:}2{:}1$ (random powder); $t=1\Rightarrow
3{:}4{:}1$ (flat texture, foil $\perp$ beam); $t=0\Rightarrow 3{:}0{:}1$ (easy axis
$\parallel$ beam). In this mode $t$ is the only free intensity parameter;
$i_1,i_2,i_3$ are derived from \eqref{eq:texture}.

\begin{figure}[htbp]\centering
  \includegraphics[width=0.98\textwidth]{apx_texture}
  \caption{Sextet texture. \emph{Left}: the intensity of the central lines $2$ and
  $5$, $i_2=4t/(2-t)$ with $t=\sin^2\theta$ \eqref{eq:texture}, varies between $0$
  (easy axis parallel to the beam, $3{:}0{:}1$) and $4$ (perpendicular foil,
  $3{:}4{:}1$), passing through the random powder $3{:}2{:}1$. \emph{Right}: the
  three situations on the same sextet; lines $2$ and $5$ grow with $t$ while the
  others do not change.}
  \label{fig:apx_texture}
\end{figure}

\subsection{Doublet and singlet}

\begin{align}
  A_{\text{dob}}(v) &= d\bigl[\,i_1\,P(v;\delta-\tfrac{\dEQ}{2},\gamma_1) + i_1 i_2\,P(v;\delta+\tfrac{\dEQ}{2},\gamma_1\gamma_2)\,\bigr],\\
  A_{\text{sing}}(v) &= d\,i_1\,P(v;\delta,\gamma_1).
\end{align}

%======================================================================
\section{Discrete fit}
%======================================================================

\subsection{Parametrization and constraints}\label{sec:disc-param}

The vector of free parameters $\bm x\in\mathbb R^{p}$ is composed of the globals
$\{b,s\}$ (plus optionally $v_{\max}$ and the fold center $c_{\text{fold}}$) and,
for each active component, a subset of
$\bm\theta_c=(\delta,\dEQ,\BHF,\gamma_1,\gamma_2,\gamma_3,d,i_1,i_2,i_3,t,\beta)$
according to its type and mode:
\begin{itemize}[leftmargin=*]
  \item \emph{Type}: a singlet uses $\{\delta,\gamma_1,d,i_1\}$; a doublet adds
        $\{\dEQ,\gamma_2,i_2\}$; a sextet uses all of them.
  \item \emph{Intensities}: in texture mode, $\{i_1,i_2,i_3\}$ are replaced by the
        single parameter $t$.
  \item \emph{Quadrupole}: in Kündig mode with fixed $\beta$, $\beta$ enters as a
        free parameter; in polycrystal or first-order mode, it does not.
  \item \emph{Fixed}: flagged parameters are removed from the active vector.
  \item \emph{Linearly constrained} $k_t=\alpha\,k_s+\beta_0$: applied at each
        evaluation (up to six iterations for chains of dependencies).
\end{itemize}
Each parameter has hard bounds $[\ell_k,u_k]$ (Appendix~\ref{ap:bounds}).

\subsection{Data statistics}\label{sec:weighting}

\subsubsection{Gauss mode}
Folding averages two independent Poisson channels; for counts $c_i$ per folded
channel, $\mathrm{Var}\approx c_i/2$, from which the $1\sigma$ uncertainty of the
normalized datum is:
\begin{equation}
  \sigma_i \;=\; \frac{\max\!\bigl(\sqrt{c_i/2},\,1\bigr)}{f_{\text{norm}}},
  \qquad \sigma_i\leftarrow\max(\sigma_i,10^{-9}).
  \label{eq:sigma-gauss}
\end{equation}

\subsubsection{Poisson mode (IRLS)}
The Poisson likelihood is approximated by iteratively reweighted least squares: the
uncertainty is evaluated on the \emph{model} (predicted counts
$\hat c_i = T_i f_{\text{norm}}$) and recomputed at each evaluation of the
optimizer, which is equivalent to a Newton step on the Poisson $-2\log\mathcal L$:
\begin{equation}
  \sigma_i^{\text{Poiss}} \;=\; \frac{\sqrt{\max(\hat c_i/2,\,1)}}{f_{\text{norm}}}
  \;=\; \frac{\sqrt{\max(T_i f_{\text{norm}}/2,\,1)}}{f_{\text{norm}}}.
  \label{eq:sigma-poisson}
\end{equation}

\subsubsection{Propagation of the calibration uncertainty}\label{sec:calib}
If the velocity calibration contributes an uncertainty $\sigma_{v_{\max}}$, its
effect on the model is added in quadrature. Since $v_i=v_{\max}(2i/(N-1)-1)$, we
have $\partial v_i/\partial v_{\max}=v_i/v_{\max}$, and by the chain rule
\begin{equation}
  \sigma_i^{\text{eff}\,2} \;=\; \sigma_i^2 \;+\;
  \left(\frac{\partial T}{\partial v}\bigg|_i \frac{v_i}{v_{\max}}\,\sigma_{v_{\max}}\right)^{\!2}.
  \label{eq:calib}
\end{equation}
The term weighs more on the flanks of the peaks (where $\partial T/\partial v$ is
large), which is where a scale error biases $\delta$ and $\BHF$.

\subsection{Cost functional and robust losses}\label{sec:cost}

The weighted residual is
\begin{equation}
  r_i(\bm x) \;=\; \frac{T(v_i;\bm x)-y_i}{\sigma_i^{\text{eff}}},
  \label{eq:residuo}
\end{equation}
and the minimized functional
\begin{equation}
  \mathcal C(\bm x) \;=\; \tfrac12 \sum_{i=1}^{N}\rho\!\left(r_i(\bm x)^2\right),
  \label{eq:cost}
\end{equation}
with $\rho$ the loss function: \emph{linear} $\rho(z)=z$ (least squares),
\emph{soft-}$\ell_1$ $\rho(z)=2(\sqrt{1+z}-1)$, or \emph{Huber}. The robust ones
use scale $f_{\text{scale}}=3$ (threshold $\sim3\sigma$), reducing the weight of
channels with spurious peaks. For linear $\rho$, $\mathcal C=\tfrac12\chi^2$.
Figure~\ref{fig:apx_losses} compares the losses and the weight each one assigns.

\begin{figure}[htbp]\centering
  \includegraphics[width=0.98\textwidth]{apx_losses}
  \caption{Robust losses versus least squares. \emph{Left}: the linear loss
  ($\chi^2$) grows as $r^2$, while soft-$\ell_1$ and Huber grow \emph{linearly} for
  large residuals. \emph{Right}: the effective weight (IRLS reweighting) is constant
  for the linear loss, but decays with $|r|$ in the robust ones, so that an outlier
  channel (spurious peak, cosmic ray) influences the fit much less.}
  \label{fig:apx_losses}
\end{figure}

\subsection{Optimization}\label{sec:opt}

\subsubsection{Trust region (TRF)}
The minimization uses \emph{Trust Region Reflective} over the bounds:
\begin{equation}
  \hat{\bm x} \;=\; \argmin_{\bm\ell\le\bm x\le\bm u}\; \mathcal C(\bm x),
\end{equation}
solving at each iteration the quadratic subproblem with a finite-difference
Jacobian:
\begin{equation}
  \min_{\bm d}\;\tfrac12\|J_k\bm d + \bm r_k\|^{2}
  \;\;\text{s.t.}\;\; \|\bm D_k\bm d\|\le\Delta_k,\;\; \bm\ell-\bm x_k\le\bm d\le\bm u-\bm x_k.
\end{equation}

\subsubsection{Multistart}
The user's point plus 8 Gaussian restarts are tried:
\begin{equation}
  \bm x^{(s)}=\bm x^{(0)}+\bm\xi^{(s)},\quad \xi^{(s)}_k\sim\mathcal N(0,\rho_k^2(u_k-\ell_k)^2),
\end{equation}
with $\rho_k=0{,}12$ for $\{\delta,\dEQ,\BHF,\gamma_1,d,v_{\max}\}$ and $0{,}08$ for
the rest; fixed seed. The one with the lowest cost is retained.

\begin{figure}[htbp]\centering
  \includegraphics[width=0.82\textwidth]{apx_multistart}
  \caption{Multistart. The cost may have local minima in addition to the global one.
  Starting from several initial points and keeping the one with the lowest cost
  reduces the risk of converging to a local minimum: in the example, only the
  central start reaches the global minimum.}
  \label{fig:apx_multistart}
\end{figure}

\subsubsection{Global optimization (optional)}
As an optional pre-pass, \emph{differential evolution} is run on the scalar cost
$\tfrac12\|\bm r\|^2$ (Sobol initialization, no polishing), and its best point is
added as a seed for the TRF polishing. It covers local minima due to label
permutation and $\BHF^{(1)}\!\leftrightarrow\!\BHF^{(2)}$ swaps in multi-sextet
spectra that the Gaussian start does not reach.

\subsection{Uncertainties and diagnostics}\label{sec:disc-stats}

\subsubsection{Covariance}
From the final Jacobian $J=U\Sigma V^\top$ (SVD, truncated singular values):
\begin{equation}
  \boxed{\;\widehat{\mathrm{Cov}}(\hat{\bm x}) = V\,\Sigma^{-2}V^\top\cdot\frac{2\,\mathcal C(\hat{\bm x})}{N-p},
  \qquad \hat\sigma_{x_k}=\sqrt{\widehat{\mathrm{Cov}}_{kk}}.\;}
  \label{eq:cov}
\end{equation}
Correlations $|\rho_{kl}|\ge0{,}95$ are reported as a warning of degeneracy.

\begin{figure}[htbp]\centering
  \includegraphics[width=0.82\textwidth]{apx_profile}
  \caption{Covariance versus profiled likelihood. The covariance \eqref{eq:cov}
  approximates the $\chi^2$ well by a parabola and gives a \emph{symmetric} error.
  The profiled likelihood scans the parameter while refitting the others and
  measures the real $\chi^2$; the cut $\Delta\chi^2=1$ gives an \emph{asymmetric}
  confidence interval, more realistic when the well is not parabolic.}
  \label{fig:apx_profile}
\end{figure}

\subsubsection{Goodness-of-fit statistics}
\begin{align}
  \chi^2 &= \sum_i r_i^2, & \nu&=N-p, & \chi^2_\nu&=\chi^2/\nu,\\
  \mathrm{AIC}&=N\ln(\mathrm{RSS}/N)+2p, & \mathrm{BIC}&=N\ln(\mathrm{RSS}/N)+p\ln N.
\end{align}

\subsubsection{Residual diagnostics}
To detect non-random structure (with $\tilde r=r-\bar r$):
\begin{align}
  \text{lag-1 autocorrelation:}\quad &\rho_1=\frac{\sum_i\tilde r_i\tilde r_{i+1}}{\sum_i\tilde r_i^2},\\
  \text{runs (Wald–Wolfowitz):}\quad &Z=\frac{R-\mu_R}{\sigma_R},\\
  \text{antisymmetric correlation:}\quad &c_{\text{as}}=\frac{\tilde{\bm r}\cdot(-\tilde{\bm r}^{\text{rev}})}{\tilde{\bm r}\cdot\tilde{\bm r}},
\end{align}
with warnings if $|\rho_1|>0{,}35$, $|Z|>2$ or $c_{\text{as}}>0{,}45$ (this last one
detects centering/calibration errors).

\subsubsection{Bootstrap}
Monte Carlo estimation of errors: $M$ replicas of synthetic data
$\tilde y^{(m)}=T(\cdot;\hat{\bm x})+\eta^{(m)}$, with $\eta^{(m)}$ Gaussian of
standard deviation $\sigma_i$ (or actual Poisson resampling
$\tilde c\sim\mathrm{Poisson}(\hat c)$ in Poisson mode); each replica is refitted
and $\hat\sigma_{x_k}^{\text{boot}}=\mathrm{std}_m(\hat x_k^{(m)})$.

%======================================================================
\section{Distribution fit}
%======================================================================

\subsection{Discretization and kernel}\label{sec:dist-kernel}

The spectrum is modeled as a continuous density $P(\BHF)$ (or $P(\dEQ)$)
discretized into $n$ \emph{bins} with equispaced centers
$c_k\in[p_{\min},p_{\max}]$. The \emph{kernel}
\begin{equation}
  K_{ik} \;=\; A_{\text{sext}}^{(d=1)}\bigl(v_i;\;\BHF=c_k\bigr)
\end{equation}
is the absorption at $v_i$ of a sextet of unit depth with field $c_k$ (variable
$\BHF$) or with $\dEQ=c_k$ and fixed $\BHF$ (variable $\dEQ$). The kernel
intensities follow $3{:}2{:}1$ by construction.

\subsection{Augmented linear system}\label{sec:dist-system}

The forward model is
\begin{equation}
  \hat y_i \;=\; b + s\,v_i \;-\; \sum_{k}K_{ik}\,P_k \;-\;\sum_{m}K^{\text{sharp}}_{im}q_m,
\end{equation}
with $P_k\ge0$ the distribution weights and $q_m\ge0$ amplitudes of optional
``sharp'' sextets. Defining the design matrix
\begin{equation}
  X = \bigl[\;\bm 1\;\big|\;\bm v\;\big|\;{-}K\;\big|\;{-}K^{\text{sharp}}\;\bigr],
  \qquad \bm z=(b,s,\bm P^\top,\bm q^\top)^\top,
\end{equation}
and the weights $W=\diag(1/\sigma_i^2)$, the problem is
\begin{equation}
  \boxed{\;
  \min_{\bm z}\;\;\bigl\|W^{1/2}(X\bm z-\bm y)\bigr\|_2^2 \;+\; \alpha\,\mathcal R(\bm P)
  \quad\text{s.t.}\quad b\ge0,\;\bm P\ge0,\;\bm q\ge0.\;}
  \label{eq:dist-min}
\end{equation}
The slope $s$ has no bound. It is solved with bounded least squares
(\code{lsq\_linear}, reflective Newton with LSMR).

\subsubsection{Column-scale preconditioning}\label{sec:cond}
For small $\BHF$ the sextet collapses and the columns of $K$ become almost
collinear with widely disparate norms, which ill-conditions \eqref{eq:dist-min}.
Each distribution column is rescaled $K_k\to K_k/s_k$ with
$s_k=\lVert W^{1/2}K_k\rVert$ and $\tilde P_k=s_k P_k$; the penalizer and the
weights are adjusted accordingly and the scale is undone at the end. It is an
\emph{exact} reparametrization (solution and degrees of freedom invariant) that
only improves numerical stability.

\subsection{Regularization}\label{sec:reg}

The inverse problem is \emph{ill-conditioned}: many different $\bm P$ reproduce the
spectrum almost equally well. Adding a term $\alpha\,\mathcal R(\bm P)$ selects a
stable solution. The value of $\alpha$ is the trade-off between fidelity and
smoothness; figure~\ref{fig:apx_regularization} shows the three regimes.

\begin{figure}[htbp]\centering
  \includegraphics[width=0.98\textwidth]{apx_regularization}
  \caption{Effect of the regularization parameter $\alpha$ on the distribution
  recovered from the same spectrum. With $\alpha$ too small (under-regularized),
  \emph{spurious oscillations} appear that fit the noise; with $\alpha$ too large
  (over-regularized) the real structure is \emph{washed out}; the value at the
  corner of the L-curve (see the distributions chapter of the user manual)
  maintains the balance. In gray, the true distribution.}
  \label{fig:apx_regularization}
\end{figure}

\subsubsection{Tikhonov (smoothness, L2)}
$\mathcal R(\bm P)=\lVert L\,\bm P\rVert_2^2$ with $L$ the \emph{second-difference}
operator $(n-2)\times n$:
\begin{equation}
  (L\bm P)_k = P_{k-1}-2P_k+P_{k+1},
\end{equation}
which penalizes curvature and produces smooth distributions. The augmented system
stacks $\bigl[\begin{smallmatrix}W^{1/2}X\\ \sqrt\alpha\,L_{\text{ext}}\end{smallmatrix}\bigr]\bm z = \bigl[\begin{smallmatrix}W^{1/2}\bm y\\ \bm 0\end{smallmatrix}\bigr]$.

\subsubsection{Total variation (edges, L1)}
$\mathcal R(\bm P)=\lVert D\,\bm P\rVert_1$ with $D$ the \emph{first difference}
$(D\bm P)_k=P_{k+1}-P_k$. It favors piecewise-constant solutions that preserve
edges (less blurring of peaks). It is solved by iteratively reweighted least
squares (IRLS): starting from the majorization
$|t|\le t^2/(2|t^{\text{prev}}|)+\text{const}$, each iteration solves a weighted
Tikhonov with
\begin{equation}
  w_k = \frac{1}{\sqrt{(D\bm P^{\text{prev}})_k^2+\varepsilon^2}},
  \qquad \mathcal R\approx \sum_k w_k\,(D\bm P)_k^2,
\end{equation}
reusing \code{lsq\_linear}. The effective operator $\sqrt{w}\,D$ is also used for
the degrees of freedom and the error bars.

\subsubsection{Maximum entropy (MaxEnt)}
Third option of \code{reg\_mode}. Instead of penalizing curvature or edges, it
maximizes the \emph{entropy} of the distribution (the most "uniform" solution
compatible with the data, without introducing unjustified structure):
\begin{equation}
  \min_{\bm z}\;\;\tfrac12\bigl\lVert (\bm y-X\bm z)/\bm\sigma\bigr\rVert_2^2
  \;-\; \alpha\,S(\bm P),\qquad
  S(\bm P) = -\sum_k P_k\,\log\!\frac{P_k}{m_k},
\end{equation}
with $m_k$ a default model (uniform). Since the entropy is \textbf{not quadratic},
\code{lsq\_linear} cannot be used: it is optimized with \textbf{L-BFGS-B} and an
analytical gradient over the physical variables $\bm z=[b,s,\bm P,\bm q]$
(\code{\_solve\_maxent} in \code{mossbauer\_distribution.py}). The effective degrees
of freedom use the Hessian corresponding to this functional.

\subsection{Poisson likelihood in the distribution}
In Poisson mode the histogram fit is repeated with IRLS reweighting of $\sigma$ from
the model (\S\ref{sec:weighting}), closing the loop through \code{lsq\_linear} until
convergence.

\subsection{Effective degrees of freedom}\label{sec:effdof}

Regularization reduces the degrees of freedom below $n$. The effective number is the
trace of the hat matrix:
\begin{equation}
  p_{\text{eff}} \;=\; \tr A(\alpha),\qquad
  A(\alpha)=X\bigl(X^\top W X+\alpha L^\top L\bigr)^{-1}X^\top W,
  \label{eq:effdof}
\end{equation}
computable as $\tr\!\bigl[(X^\top WX+\alpha L^\top L)^{-1}X^\top WX\bigr]$. It is used
in $\chi^2_\nu=\chi^2/(N-p_{\text{eff}})$, avoiding the bias of counting the $n$ bins
as independent parameters. It is invariant under the column rescaling of
\S\ref{sec:cond}.

\subsection{Selection of the regularization parameter}\label{sec:alpha}

\subsubsection{Dimensionless $\alpha$}\label{sec:alpha-dimless}
The $\alpha$ set by the user is \textbf{not} the raw weight of \eqref{eq:dist-min}.
The data term is a $\chi^2\sim N$, but the penalizer is in units of absorption;
without normalization, the "useful" $\alpha$ would depend on $N$, on the noise
$\sigma$ and on the absorption depth, and the same number would mean different
things in each spectrum. That is why the interface's $\alpha$ is
\textbf{made dimensionless} before applying it:
\begin{equation}
  \alpha_{\text{eff}} \;=\; \alpha\;\lambda_{\text{ref}}\;c_0,
  \qquad
  \lambda_{\text{ref}} = \frac{\lVert A_{\text{dist}}\rVert_F^2}{\lVert L_{\text{scaled}}\rVert_F^2},
  \qquad c_0 = 10^{-4},
\end{equation}
where $A_{\text{dist}}$ and $L_{\text{scaled}}$ are the distribution part of
$W^{1/2}X$ and the \emph{whitened} penalizer (after the column rescaling of
\S\ref{sec:cond}), $\lVert\cdot\rVert_F$ the Frobenius norm and $c_0=$
\code{ALPHA\_REF\_SCALE}. The augmented system uses $\sqrt{\alpha_{\text{eff}}}\,L$
instead of $\sqrt\alpha\,L$; the same for $p_{\text{eff}}$ \eqref{eq:effdof} and the
error bars. The effect: $\log_{10}\alpha\approx 0$ is the \textbf{natural
balance} data/smoothness for \emph{any} spectrum, and the elbow of the L-curve
stays stable against $N$, noise and depth (typical useful values
$-1\ldots+2$). In 2D the same is done for each axis ($\alpha_B$, $\alpha_Q$) with
its own $\lambda_{\text{ref}}$.

\medskip
Three criteria are computed over a logarithmic grid of $\alpha$.

\paragraph{L-curve (maximum curvature).} With
$\rho(\alpha)=\log\lVert W^{1/2}(X\hat z-y)\rVert$ and
$\eta(\alpha)=\log\lVert L\hat{\bm P}\rVert$, the corner maximizes
\begin{equation}
  \kappa(\alpha)=\frac{\rho'\eta''-\eta'\rho''}{\bigl((\rho')^2+(\eta')^2\bigr)^{3/2}}.
\end{equation}

\paragraph{Trade-off.} Minimum normalized distance to the origin in the
$(\eta,\rho)$ plane.

\paragraph{Generalized cross-validation (GCV).}
\begin{equation}
  \mathrm{GCV}(\alpha)=\frac{\lVert W^{1/2}(\bm y - X\hat z)\rVert^2}{\bigl(N-p_{\text{eff}}(\alpha)\bigr)^2},
  \qquad \hat\alpha=\argmin_\alpha \mathrm{GCV}(\alpha),
\end{equation}
which does not require knowing the noise scale and is optimal in prediction loss.

\subsection{Error bars of \texorpdfstring{$P(\BHF)$}{P(BHF)}}\label{sec:perr}

For the regularized estimator $\hat{\bm z}=H^{-1}X^\top W\bm y$ with
$H=X^\top WX+\alpha L^\top L$, the linearized covariance is
\begin{equation}
  \widehat{\mathrm{Cov}}(\hat{\bm z}) \;=\; H^{-1}\,(X^\top WX)\,H^{-1}\cdot\frac{\chi^2}{N-p_{\text{eff}}},
\end{equation}
scaled by the effective $\chi^2_\nu$ (same convention as the discrete fit). The
square root of the diagonal corresponding to the bins gives $\sigma_{P_k}$. The band
$P_k\pm\sigma_{P_k}$ is clipped to $P\ge0$ when plotted (non-negativity makes it
asymmetric in bins pinned to zero). In TV mode the effective operator $\sqrt{w}\,D$
is used instead of $L$.

\subsection{Parametric distributions}
As an alternative to the regularized histogram, closed forms with few parameters are
fitted (nonlinear least squares, no $\alpha$):
\begin{align}
  \text{Gaussian:}\quad & P_k\propto \exp\!\bigl(-\tfrac12((c_k-\mu)/\varsigma)^2\bigr), & (b,s,A,\mu,\log\varsigma),\\
  \text{VBF:}\quad & P_k = \sum_{j=1}^{N} A_j\,\mathcal N(c_k;\mu_j,\sigma_j), & (b,s,\{A_j,\mu_j,\log\sigma_j\}_{j=1}^N),\\
  \text{Binomial:}\quad & P_k\propto \binom{n-1}{k-1}p^{k-1}(1-p)^{n-k}, & (b,s,A,\mathrm{logit}\,p),\\
  \text{Fixed:}\quad & P_k=\phi_k\ \text{(preloaded)}, & (b,s,A).
\end{align}

\paragraph{VBF (Voigt-Based Fitting, Rancourt--Ping 1991).} It generalizes the
Gaussian to a sum of $N$ Gaussians ($N=1\ldots6$) over a Voigt kernel, each one
with amplitude $A_j$, mean $\mu_j$ and width $\sigma_j$ (stored ordered by $\mu$).
Since the Gaussian kernel $\mathcal N$ is \textbf{normalized to unit area}, the
coefficient $A_j$ \emph{is} the \textbf{area} of that component over the grid; its
fraction $A_j/\sum_i A_i$ gives the \textbf{population} of each environment, which
Fitbauer reports as "area \%". The \emph{Gaussian} form is simply the case $N=1$.
Appendix~\ref{ap:vbf} develops the complete theory of the method (forward model,
Voigt profile, parametrization and best practices).

\subsection{Correlation \texorpdfstring{$\delta(H)$}{d(H)} and \texorpdfstring{$\dEQ(H)$}{DEQ(H)} in the kernel}\label{sec:dist-corr}
By default all the subspectra of the kernel share the same $\delta$ and $\dEQ$, and
only the field varies. Optionally $\delta$ and $\dEQ$ are coupled to the bin field
through two slopes (Le~Caër--Dubois / Wivel--M\o{}rup model):
\begin{equation}
  \delta_k = \delta_0 + \kappa_\delta\,(c_k-H_{\text{ref}}),\qquad
  \Delta E_{Q,k} = \Delta E_{Q,0} + \kappa_q\,(c_k-H_{\text{ref}}),
\end{equation}
with $H_{\text{ref}}$ a reference field. Thus the kernel $K_{ik}$ is built using
$(\delta_k,\Delta E_{Q,k})$ in each column instead of global values. With
$\kappa_\delta=\kappa_q=0$ the classic case is recovered. The slopes are refined in
the outer loop (\S\ref{sec:refine}) and apply to the Histogram and to the VBF.

\subsection{Global refinement (VARPRO)}\label{sec:refine}
When enabled, a nonlinear outer loop optimizes the shape parameters
$\bm\eta=(\delta_g,\log\gamma_g,\dots)$; at each evaluation the inner linear problem
\eqref{eq:dist-min} is solved \emph{in full}, returning the residual to the
optimizer (separable VARPRO-type scheme):
\begin{equation}
  \min_{\bm\eta}\;\bigl\|W^{1/2}\bigl[X(\bm\eta)\,\hat{\bm z}(\bm\eta)-\bm y\bigr]\bigr\|^2,
  \qquad \hat{\bm z}(\bm\eta)=\argmin_{\bm z}\eqref{eq:dist-min}.
\end{equation}

%======================================================================

\section{Appendices}

\subsection{Sextet calibration}\label{ap:calib}
The reference positions at $B_0=33{,}0\,$T are the \textbf{published} ones of
$\alpha$-Fe (standard velocity reference), with no additional rescaling factor:
\begin{equation}
  r^{(0)}=\tfrac12[-10.657,-6.167,-1.677,1.677,6.167,10.657]
        = [\mp5{,}3285,\ \mp3{,}0835,\ \mp0{,}8385]\ \text{mm/s}.
\end{equation}
The positions for a field $\BHF$ scale linearly, $r^{(0)}\cdot(\BHF/33{,}0)$. From
them the Zeeman splittings used in \eqref{eq:Hkundig} are derived:
\begin{equation}
  \omega_e^{(0)} = r^{(0)}_2-r^{(0)}_1 \approx 2{,}245\ \text{mm/s},\qquad
  \omega_g^{(0)} = -\bigl[(r^{(0)}_4-r^{(0)}_3)+\omega_e^{(0)}\bigr] \approx -3{,}922\ \text{mm/s}.
\end{equation}
\begin{nota}[No factor 32.95]
Previous versions of this document included a factor $(B_0/32{,}95)$ that
over-scaled the positions by $\sim0{,}15\%$. The code uses the published positions
\emph{as is} (\code{LINE\_POS\_33T} in \code{core/constants.py}), so that factor has
been removed (and with it $\omega_e^{(0)}$, $\omega_g^{(0)}$ go from
$2{,}248/-3{,}928$ to $2{,}245/-3{,}922$).
\end{nota}

\subsection{Parameter bounds}\label{ap:bounds}
\begin{center}
\begin{tabular}{@{}lll@{}}
\toprule
Parameter & Bound & Unit\\
\midrule
$b$ (baseline) & $[0{,}7,\ 1{,}3]$ & --\\
$s$ (slope) & $[-5\cdot10^{-3},\ 5\cdot10^{-3}]$ & mm$^{-1}$s\\
$\delta$ & $[-2,\ 3]$ & mm/s\\
$\dEQ$ & $[-4,\ 4]$ & mm/s\\
$\BHF$ & $[0,\ 60]$ & T\\
$\gamma_1$ (FWHM) & $[0{,}06,\ 4]$ & mm/s\\
$\gamma_2,\gamma_3$ & $[0{,}2,\ 6]$ & (relative)\\
$d$ (depth) & $[0,\ 0{,}30]$ & --\\
$i_1,i_2,i_3$ & $[0,9]/[0,6]/[0,3]$ & (relative)\\
$t$ (texture) & $[0,\ 1]$ & --\\
$\beta$ & $[0,\ 90]$ & degrees\\
\bottomrule
\end{tabular}
\end{center}

\subsection{Faddeeva function}
$w(z)=e^{-z^2}\erfc(-iz)$, evaluated by \code{scipy.special.wofz}. The real part
gives the Voigt profile \eqref{eq:voigt}; for $z=iy$ with $y\in\mathbb R$,
$w(iy)=e^{y^2}\erfc(y)$ is real, which allows the analytical normalization to the
peak.

\subsection{Summary of the optimization stack}
\begin{center}
\begin{tabular}{@{}lll@{}}
\toprule
\textbf{Mode} & \textbf{Functional} & \textbf{Solver}\\
\midrule
Discrete & $\tfrac12\sum_i\rho(r_i^2)$ & \code{least\_squares} (TRF) [+ optional DE]\\
Tikhonov distribution & $\|W^{1/2}(X\bm z-\bm y)\|^2+\alpha\|L\bm P\|^2$ & \code{lsq\_linear} (bounded)\\
TV distribution & $\|W^{1/2}(X\bm z-\bm y)\|^2+\alpha\|D\bm P\|_1$ & IRLS over \code{lsq\_linear}\\
Parametric distribution & $\|W^{1/2}(X(\bm\theta)\bm z-\bm y)\|^2$ & \code{least\_squares} (NL)\\
Global refinement & $\|W^{1/2}(X(\bm\eta)\hat{\bm z}(\bm\eta)-\bm y)\|^2$ & outer NL + inner bounded\\
\bottomrule
\end{tabular}
\end{center}
